\cvsection{Expériences professionnelles}

\begin{cventries}
  \cventry
    {Stagiaire ingénieur de recherche en IA}
    {INRIA, équipe Mnémosyne}
    {Bordeaux, France}
    {Mars 2026 -- Septembre 2026}
    {
      \begin{cvitems}
        \item {Conception et développement en autonomie, avec le suivi de mes encadrants, d’un agent IA multi-étapes répondant au besoin d’aider les chercheurs à trouver des articles scientifiques.}
        \item {Construction de l’agent sans framework d’orchestration, intégration de l’API OpenAI et conception de ses outils afin de maîtriser l’exécution, l’état, la mémoire et le contexte.}
        \item {Développement du backend en Python/FastAPI avec traitements asynchrones, API REST documentées avec OpenAPI et persistance MongoDB.}
        \item {Mise en œuvre d’un pipeline RAG avec chunking, embeddings, recherche sémantique et ChromaDB.}
        \item {Développement d’une interface React/Next.js, d’un outil Python de collecte de métadonnées scientifiques et de webhooks transmettant les résultats de tâches planifiées.}
        \item {Conteneurisation de composants avec Docker, mise en place d’un pipeline CI/CD sur GitHub et évaluation d’OpenClaw dans un environnement isolé avec une attention portée à la sécurité.}
        \item {Gestion d’une application web conteneurisée permettant d’utiliser des LLM installés localement sur un Mac Studio, incluant l’installation et la configuration des modèles.}
        \item {Comparaison de Qwen, Gemma et GLM et de plusieurs quantifications sur une infrastructure HPC avec Slurm selon la latence, la TTFT, la qualité et la taille du contexte, puis contribution à un article scientifique.}
      \end{cvitems}
    }

  \cventry
    {Stagiaire de recherche en apprentissage par renforcement}
    {Université Karunya}
    {Coimbatore, Tamil Nadu, Inde}
    {Juin 2025 -- Sept. 2025}
    {
      \begin{cvitems}
        \item {Conception, développement et comparaison d’une planification de trajectoire déterministe et d’une approche fondée sur l’apprentissage par renforcement pour la navigation autonome d’un robot hospitalier.}
        \item {Déploiement et optimisation du logiciel et du modèle d’apprentissage par renforcement sur un Raspberry Pi embarqué.}
      \end{cvitems}
    }
\end{cventries}
